% ============================================================
%  2. PRESENTACIÓN DEL PROBLEMA Y DE LA NECESIDAD  (Subdocumento 2)
%  Informe 1 — Presentación del problema y de la necesidad
%  Fuente: Informe (1).md (pp. 92-400); Caso 02 — Distribuidora Puelche
% ============================================================
\chapter{Presentación del Problema y de la Necesidad}
\label{sec:problema}

\section{Contextualización del Entorno}

La \textbf{Distribuidora Puelche S.A.} es una empresa familiar fundada en 1974, hoy
dirigida por su tercera generación. Su especialidad es la distribución mayorista de
consumo masivo (alimentos, bebidas, aseo del hogar y cuidado personal) en el centro-sur
de Chile, desde la Región de O'Higgins hasta la Región de La Araucanía. La red comprende
seis instalaciones distribuidas en cuatro regiones (Tabla 14.1 y RT-21.16 del Caso).

La empresa opera desde dos centros de distribución principales: el primero en Talca
(18.000\,m²) y el segundo en Concepción (9.000\,m²), más tres plataformas de
\emph{cross-docking} en Curicó, Chillán y Los Ángeles. Mantiene de forma activa
\textbf{14.200 clientes}, distribuidos en su mayoría entre almacenes y minimarkets del
canal tradicional (11.600 clientes), restaurantes y casinos del segmento \emph{food
service} (2.100) y cadenas de supermercados regionales (hasta 500 puntos de entrega en
este sector).

Solo en 2025, la compañía procesó 31.000 pedidos mensuales, movilizó 2,4 millones de
unidades y emitió cerca de 34.000 documentos tributarios, con ventas anuales de \$148.000
millones. Su dotación propia alcanza las 640 personas, complementada por
aproximadamente 160 conductores pertenecientes a 10 empresas transportistas externas.

\subsection{El detonante: el retiro sanitario de marzo de 2026}

El 3 de marzo de 2026, el gerente de aseguramiento de calidad de uno de los principales
proveedores de lácteos del país notificó a Puelche la necesidad de ejecutar un retiro
preventivo por una desviación detectada en un lote identificado como \textbf{24-0217} de
queso fresco. La respuesta esperada era identificar, en el menor tiempo posible, a cuáles
clientes había llegado dicho lote, la cantidad y la documentación que traía. La jefa
tardó nueve días en responder y, para peor, la respuesta no fue exacta, sino una
estimación ni siquiera certera. Al ser prácticamente imposible conocer los puntos de
entrega afectados, la empresa debió retirar de forma total el producto sobre los 14.200
clientes.

\begin{itemize}
  \item Pérdida financiera: \$31.000.000 de pérdidas directas producto del retiro.
  \item 4.800 kilogramos de queso fresco retirados del mercado.
  \item Un sumario sanitario abierto por la autoridad.
  \item Suspensión como distribuidor autorizado por el proveedor durante seis meses,
        equivalente al 6\,\% de las ventas totales.
\end{itemize}

Este evento expuso una realidad que se intuía pero no se había cuantificado: \textbf{no
existe trazabilidad}. El registro de lote de un producto, en el mejor de los casos, solo
lleva una planilla impresa archivada en una carpeta y, según la propia distribuidora,
alrededor del \textbf{41\,\% de los casos no existe en absoluto} esa planilla. No hay
ningún sistema operativo que permita responder en tiempo real adónde fue a parar un
producto específico.

\subsection{Las presiones externas que transforman la crisis en amenaza estratégica}

El retiro sanitario no fue el único problema del comité ejecutivo en el primer trimestre
de 2026. De forma paralela se crearon dos presiones externas:

\begin{itemize}
  \item \textbf{Irrupción de un competidor nacional} que abrió operación en Talca ocho
        meses antes. Ofrece entregas en 24 horas y una aplicación móvil donde el
        almacenero realiza su pedido de forma autónoma, consulta el stock disponible y
        sabe la llegada aproximada del camión. Puelche, en contraste, entrega entre 48 y
        72 horas y toma el pedido mediante un preventista que visita al cliente una vez
        por semana.
  \item \textbf{El cliente más importante} (una cadena de supermercados que representa el
        11\,\% de la venta) notificó sus condiciones comerciales para el período que
        comienza en enero de 2029:
        \begin{enumerate}
          \item Pedidos exclusivamente por vía electrónica estructurada.
          \item Aviso de despacho anticipado.
          \item Ventana de entrega de treinta minutos con penalización por
                incumplimiento.
          \item Prueba de entrega digital obligatoria.
        \end{enumerate}
        Hoy Puelche no cumple ninguna de estas condiciones: la ventana de entrega es de
        10 horas (8:00 a 18:00), los pedidos no se reciben por vía electrónica y la
        prueba de entrega es una guía de papel firmada a mano.
\end{itemize}

\subsection{Metas convocantes de la licitación}

Estas tres presiones llevaron al comité ejecutivo a convocar una licitación para
desarrollar e implementar una solución integral de gestión logística, con las siguientes
metas:

\begin{enumerate}
  \item Entregas completas y a tiempo: de 82,4\,\% a más del 95\,\%.
  \item Cumplimiento de lo pedido: de 91,3\,\% a más del 97\,\%.
  \item Plazos de entrega: de 48--72 horas a 24 horas.
  \item Ventana de entrega comprometida: disminuida hasta 30 minutos para 2029.
  \item Tiempo de respuesta ante un retiro sanitario: de horas, con evidencia exacta.
  \item Registro continuo y automático de temperatura en cadena de frío.
  \item Costo de servir por cliente calculado por entrega.
  \item Pedidos recibidos mediante vía electrónica estructurada.
\end{enumerate}

\section{Marco Legal y Regulatorio}

El diseño, la construcción y la operación de la plataforma digital de misión crítica se
encuentran condicionados por un entramado normativo cuyo cumplimiento no constituye una
formalidad declarativa, sino un conjunto de \textbf{restricciones técnicas} que inciden
directamente en la arquitectura, los controles de seguridad, la estructura de datos y los
flujos de integración.

\subsection{Marco legal nacional}

\begin{table}[htbp]
\centering
\small
\caption{Marco legal nacional y su impacto técnico}
\label{tab:marco-legal}
\begin{tabularx}{\textwidth}{@{}L{4.4cm}X@{}}
\toprule
\textbf{Cuerpo normativo} & \textbf{Impacto técnico en la solución} \\
\midrule
Ley N° 21.719 y Ley N° 19.628 (protección de datos personales) &
Modelo de acceso RBAC + ABAC, cifrado a nivel de campo AES-256 para datos sensibles
(RT-11.10), registro de actividades de tratamiento, anonimización/seudonimización en
ambientes no productivos (RT-11.25) y residencia de datos. \\
\addlinespace
Ley N° 21.663 (Ley Marco de Ciberseguridad e Infraestructura Crítica) &
Modelo Zero Trust (NIST SP 800-207), SOC 24$\times$7, notificación de brechas en 24
horas, segmentación de redes, sin acceso directo a datos desde Internet. \\
\addlinespace
Ley N° 21.459 (Delitos Informáticos) &
Pistas de auditoría inalterables (RT-16.07), registro de quién/dispositivo/fecha/valores
antes-después (RT-16.06), WAF y protección DDoS capas 3--4--7 (RT-11.07), correlación
SIEM (RT-11.15). \\
\addlinespace
Ley N° 19.799 (Documentos electrónicos y firma electrónica) &
Prueba de entrega digital articulada con la guía de despacho electrónica y su acuse de
recibo (RT-16.18). \\
\addlinespace
Ley N° 17.336 (Propiedad intelectual) &
Proceso formal de aprobación de dependencias de terceros (RT-11.26), inventario de
componentes SBOM (RT-11.23), SLSA nivel 3 (RT-11.24), licencias a nombre del CLIENTE. \\
\addlinespace
Ley N° 20.393 y Ley N° 21.595 (responsabilidad penal de las personas jurídicas) &
Aprobación dual de transacciones de alto riesgo (RT-16.03), elevación de privilegios
temporal y grabada (RT-12.06), rendición de efectivo registrada en el momento del cobro. \\
\addlinespace
Reglamento Sanitario de los Alimentos (RSA, D.S. 977) &
Trazabilidad lote-a-lote desde recepción hasta punto de entrega, estándares GS1 (GTIN,
SSCC, GLN), captura automatizada de temperatura con sensores IoT en cámaras y vehículos. \\
\addlinespace
Normativa tributaria (DTE) &
La solución no emite documentos tributarios (permanece en el ERP) pero se integra
bidireccionalmente con él, respetando Restricción no negociable N° 4 (sin ``dos
verdades''). \\
\bottomrule
\end{tabularx}
\end{table}

\subsection{Estándares internacionales y marcos de referencia}

La solución se adhiere a estándares internacionales que estructuran sus decisiones de
arquitectura:

\begin{itemize}
  \item \textbf{ISO/IEC 27001/27002:} SGSI con matriz de controles trazable (RT-11.05);
        certificación vigente o plan de certificación en los primeros 12 meses.
  \item \textbf{ISO/IEC 27017/27018:} seguridad y datos personales en nube (RT-11.06).
  \item \textbf{ISO/IEC 27005:} gestión de riesgos de seguridad de la información
        (RT-11.02).
  \item \textbf{ISO 22301 e ISO/IEC 27031:} continuidad de negocio y continuidad TIC,
        con RTO máximo de 4 horas y RPO máximo de 15 minutos (RT-07.04).
  \item \textbf{Disponibilidad 99,95\,\%} del Capítulo 6 de las Bases Técnicas
        Transversales (RT-06.01), alineado con instalaciones Tier III del Uptime
        Institute.
  \item \textbf{NIST CSF 2.0 y NIST SP 800-207} (Zero Trust) como fundamento de
        seguridad (RT-11.01).
  \item \textbf{OWASP ASVS 4.0, OWASP Top 10 y OWASP API Security Top 10}, con CIS
        Benchmarks para endurecimiento on-premise (RT-03.15).
  \item \textbf{ISO/IEC 20000-1 e ITIL 4:} gestión de servicios de TI.
  \item \textbf{ISO/IEC 25010 e ISO/IEC 25012:} calidad de producto y de datos
        (RT-05.04).
  \item \textbf{WCAG 2.2 (nivel AA):} accesibilidad de todas las interfaces.
  \item \textbf{Estándares GS1:} identificación e interoperabilidad (RT-05.23).
\end{itemize}

\section{Modelo Operativo Actual (Procesos AS-IS)}

Este apartado describe el flujo operativo diario tal como se ejecuta actualmente, con
propósito diagnóstico: identificar los puntos donde la ausencia de digitalización, la
duplicación de registros y los silos de información generan ineficiencias medibles y
riesgos cuantificados.

\subsection{Abastecimiento y recepción de mercadería}

El área de compras genera las órdenes de compra en el ERP (implantado en 2017) a partir
de una sugerencia de reposición que no produce el propio sistema, sino una planilla de
cálculo mantenida por el jefe de abastecimiento, que considera la venta de las últimas
ocho semanas y un factor de ajuste manual por temporada. Las promociones no ingresan al
cálculo de forma automática.

Los 180 proveedores activos confirman las órdenes por correo electrónico; ninguno lo hace
por vía electrónica estructurada. La recepción física se realiza sin cita previa, con
filas de hasta cinco horas en temporada alta. La verificación es enteramente manual: el
recepcionista cuenta la mercadería contra la guía en papel, revisa fechas de vencimiento
por muestreo visual y anota el lote en un campo de texto libre del sistema cuando trae
la información visible y cuando efectivamente se registra. \textbf{En el retiro de marzo
de 2026, el 41\,\% de las recepciones del producto involucrado carecía de registro de
lote.}

\subsection{Almacenamiento e inventario}

Opera un sistema de gestión de almacenes (WMS) instalado en 2013, usados exclusivamente
en Talca. Concepción controla su inventario mediante planillas locales y las tres
plataformas de \emph{cross-docking} carecen por completo de control sistematizado. La
asignación de ubicaciones se realiza según el criterio personal del jefe de bodega, sin
regla documentada.

El conteo cíclico se ejecuta sobre 3.400 posiciones mensuales en el turno de noche con
planillas impresas. La diferencia promedio del conteo es del \textbf{2,3\,\%} del valor
contado, y los ajustes se ejecutan sin investigación de causa raíz.

\subsection{Preventa y toma de pedidos}

62 preventistas recorren rutas fijas con frecuencia semanal o quincenal. Operan en
terreno, de pie a la intemperie, con un dispositivo que ejecuta una aplicación adquirida
en 2016 a un proveedor que ya no existe: sin soporte, sin compatibilidad y sin
actualizaciones. La aplicación no muestra stock disponible, crédito del cliente, deuda
vencida ni historial de compra; el preventista trabaja ``a ciegas''.

\begin{itemize}
  \item Se venden productos que no existen en stock: el \textbf{7,8\,\%} de las líneas de
        pedido presentan quiebre en el momento de la entrega.
  \item Se toman pedidos a clientes con crédito bloqueado, y el pedido queda retenido sin
        notificación.
  \item En zonas rurales con tramos sin señal de hasta dos horas, la aplicación guarda
        pedidos localmente, pero su comportamiento al recuperar cobertura es
        impredecible: no envía, o envía duplicados. El equipo comercial estima que
        ``pasa todas las semanas''.
  \item La información comercial observada en terreno (vencidos, exhibidores del
        competidor, riesgo de cierre) se pierde estructuralmente.
\end{itemize}

\subsection{Planificación de rutas}

La asignación diaria de aproximadamente 1.400 entregas a los 96 camiones disponibles (42
propios y 54 de transportistas) la realiza \textbf{una única persona}: Hugo Maldonado,
con 21 años de antigüedad, que consume entre 3,5 y 4 horas diarias y se apoya en una
planilla de cálculo de 11 hojas que ningún otro integrante comprende ni replica. En esa
planilla y en su memoria reside conocimiento crítico: restricciones de carga en puentes
rurales, horarios de atención, preferencias de conductor y condiciones estacionales.
Cuando el planificador toma vacaciones, el cumplimiento de entrega cae entre 6 y 9
puntos porcentuales. \textbf{Se jubila en dos años.}

El resultado es una planificación subóptima y frágil: la ocupación promedio de los
camiones en volumen es del 68\,\%, con días en que un camión sale con apenas el 41\,\%
de su capacidad.

\subsection{Preparación de pedidos y carga}

La preparación se ejecuta en el turno de noche (22:00 a 06:00) con aproximadamente 120
personas y una rotación anual del 38\,\%. El preparador recibe una hoja de \emph{picking}
impresa y recorre el almacén con un transpaleta. Los faltantes se anotan manualmente sin
su causa, impidiendo el análisis de causa raíz. En la cámara de congelado
($-$22\,°C), los preparadores trabajan con guantes térmicos en turnos de 30 minutos; los
dispositivos convencionales no operan de forma confiable a esa temperatura y no existe
cobertura de señal en su interior (Restricción no negociable N° 7).

La carga se realiza por orden de ruta invertido, dependiendo de la secuencia que reside
exclusivamente en la planilla del planificador.

\subsection{Transporte, entrega y cobro}

Los camiones salen entre las 05:30 y las 07:00 portando guías de despacho impresas en
papel. El conductor lleva un termógrafo manual que debe leer tres veces durante el viaje;
no existe registro continuo ni automático de temperatura.

El cliente firma la guía en papel. Las devoluciones se anotan manualmente. El 38\,\% de
las ventas del canal tradicional se cobra al contado, y en rutas rurales un conductor
puede circular con \textbf{más de un millón de pesos en efectivo}, que se rinde al día
siguiente, con diferencias de rendición que promedian \$4,2 millones mensuales sin
investigación de causa.

\subsection{Devoluciones, mermas y activos retornables}

Las devoluciones y mermas se identifican sin causa, imposibilitando distinguir entre
errores de preparación, daños en tránsito o rechazos del cliente. No existe control
estructurado sobre los activos retornables (canastillos y \emph{pallets}).

\subsection{Síntesis de ineficiencias estructurales del modelo AS-IS}

El modelo actual se caracteriza por silos de información, reingreso manual de datos,
ausencia de trazabilidad de lote, operación sin conectividad sin mecanismos de
deduplicación, dependencia de conocimiento no documentado y ausencia de registro continuo
de temperatura y de costo de servir.

\diagrama{as_is_preventa_reparto.png}{Flujo AS-IS: preventa, toma de pedidos y reparto actual}{fig:as_is_preventa}

\diagrama{as_is_preparacion_despacho.png}{Flujo AS-IS: preparación nocturna y despacho (ventana 05:30--07:00)}{fig:as_is_preparacion}

\section{Diagnóstico de Ineficiencias}

\begin{table}[htbp]
\centering
\small
\caption{Síntesis diagnóstica por proceso}
\label{tab:diagnostico}
\begin{tabularx}{\textwidth}{@{}L{4.2cm}X@{}}
\toprule
\textbf{Proceso} & \textbf{Ineficiencia identificada} \\
\midrule
Preventas &
Operación ``a ciegas'' sin stock, crédito ni historial; pedidos no enviados o
duplicados; aplicación de 2016 sin soporte. \\
\addlinespace
Planificación de rutas &
Conocimiento crítico en una sola persona; ocupación vehicular de 68\,\%;
planificación de 3,5--4 horas diarias. \\
\addlinespace
Registro del pedido y protocolos &
Reingreso manual; sin deduplicación ni resolución determinista de conflictos;
bloqueos por crédito sin notificación. \\
\addlinespace
Trazabilidad de lote &
Inexistente (41\,\% de recepciones sin registro); retiro sanitario de nueve días y no
exacto. \\
\addlinespace
Cadena de frío &
Sin registro continuo ni automático; lectura manual de termógrafo tres veces por
viaje. \\
\addlinespace
Cobranza en efectivo &
Hasta \$1M en efectivo por conductor; rendición al día siguiente; diferencias de
\$4,2M mensuales sin investigación. \\
\addlinespace
Activos retornables y mermas &
Sin control estructurado ni causa raíz. \\
\addlinespace
Costo logístico &
Prorrateo por zona de 2016 sin vigencia técnica; costo de servir por entrega
inexistente. \\
\bottomrule
\end{tabularx}
\end{table}

\section{Análisis de Impacto}

\subsection{Impacto financiero directo}

El retiro sanitario de marzo de 2026 generó \$31.000.000 en pérdidas directas y una
suspensión comercial de seis meses equivalente al 6\,\% de las ventas. A ello se suman
las diferencias de rendición de efectivo de \$4,2 millones mensuales sin investigación y
la ocupación vehicular subóptima de 68\,\%.

\subsection{Impacto operacional}

La dependencia de una única persona para planificar 1.400 entregas diarias, la pérdida
de datos por doble digitación y el reingreso manual en recepción, \emph{picking} y conteo
cíclico comprometen la continuidad operacional y la confiabilidad de los registros.

\subsection{Impacto regulatorio y comercial}

El sumario sanitario, la exigencia del canal moderno para 2029 y la competencia con
entregas en 24 horas amenazan la posición competitiva. El incumplimiento de las
condiciones del cliente más importante (11\,\% de la venta) representa un riesgo
estratégico de primer orden.

\section{Problemas por Área}

Los problemas detectados se agrupan en cinco áreas que estructuran la propuesta de
solución:

\begin{enumerate}
  \item \textbf{Área 1 — Trazabilidad y calidad sanitaria:} registro de lote
        estructurado, monitoreo continuo de cadena de frío y respuesta a retiros
        sanitarios.
  \item \textbf{Área 2 — Preventa y gestión comercial en terreno:} visibilidad de stock,
        crédito e historial, deduplicación determinista y captura de inteligencia
        comercial.
  \item \textbf{Área 3 — Planificación de rutas y operación de reparto:} motor de
        ruteo, prueba de entrega digital, cobro en el momento y control de envases.
  \item \textbf{Área 4 — Gestión de bodega e inventario:} WMS unificado, \emph{picking}
        digital, conteo cíclico estructurado y registro de causas.
  \item \textbf{Área 5 — Información de gestión y costo de servir:} capa analítica,
        costeo por entrega y capa de datos para la toma de decisiones.
\end{enumerate}

\section{Actores Afectados}

El diagnóstico identifica una amplia diversidad de actores afectados, cuyo contexto
operacional (frío extremo, lluvia, sin señal, una sola mano) condiciona el diseño de la
solución. Entre los principales: preventistas (62, proyección 70), conductores propios
(42) y peonetas (42), conductores de 10 empresas transportistas, preparadores de bodega
(120, turno de noche), recepcionistas, planificador de rutas, jefe de bodega, jefe de
abastecimiento, equipo de crédito y cobranza, jefe de TI (único responsable con 4
personas), y los clientes de los canales tradicional, moderno y \emph{food service}.

\section{Registro de Supuestos del Caso}

\pendiente{} Consolidar aquí el registro de supuestos del caso, incluyendo los
supuestos por épica declarados en el documento de trabajo (entorno, datos, acuerdos con
terceros, condiciones de mercado y vigencia de las parametrizaciones). Este registro se
detalla en la pestaña de supuestos y se complementará en el Informe 2 junto con el
análisis de riesgos.

\vspace{4pt}
{\small\color{lafroxGrato}\pendiente{} Enriquecer las secciones 2.3 (procesos AS-IS) con
los diagramas de flujo correspondientes y las secciones 2.7 (actores) y 2.8 (supuestos)
con el detalle completo del documento de trabajo.}
